\begin{papercard}{zych2026}
\cf{Stages} \stagetag{A} \stagetag{D} \stagetag{R}. A secondary, exploratory analysis of an existing commercial assessment dataset: ORCHA Baseline Review results for 436 mental-health apps. It uses statistical and machine-learning methods rather than a researcher-drawn store sample. \emph{Note:} the appended copy has a repository cover page, so article page 1 is PDF page 2.

\cf{Aim and RQs} Examine quality, target users, features, evidence and data privacy of mental-health apps through seven exploratory questions, e.g.\ how apps score per domain, whether cost matters, which data are collected \pc{zych2026}{3}{How do mental health apps score on the OBR assess}.

\cf{Data source} ORCHA (a UK digital-health compliance company) provided OBR V6 assessments of 2{,}127 apps published January 2021 to January 2022. Each app was assessed by two ORCHA reviewers, with a third for disputes. Reviewer training takes about six months \pc{zych2026}{3}{The assessment and app catego}. The OBR has about 300 objective, mostly yes/no items in three domains (professional and clinical assurance, data privacy, user experience), each scored 0--100 and combined into an overall score \pc{zych2026}{4}{The OBR V6 tool is the latest version of the ORCHA}. Subset rule: only ORCHA's ``mental health app'' category, giving 436 apps \pc{zych2026}{4}{selection criterion was that an app had been categorized as}. Apps are tiered by the NICE Evidence Standards Framework (A/B/C by function, risk and regulatory status), and weighting differs by tier \pc{zych2026}{4}{has produced an ESF that can be used to classify digital}.

\cf{How ORCHA sampled} The underlying dataset follows ORCHA's routine approach: categorise apps, order categories by downloads, assess the most downloaded first, then the next, and so on. A pilot ran on 30 highly compliant apps \pc{zych2026}{15}{assessing the most downloaded digital health app in each category}.

\cf{Analysis} Descriptives with quantiles and median/IQR; Wilcoxon rank-sum with Bonferroni correction; random-forest feature importance for the ORCHA score; association-rule mining (apriori) on co-collected data types; k-modes clustering with NbClust (30 indices) on binary data-type variables \pc{zych2026}{4}{A random forest algorithm was used to determine feature}. Ethics: secondary-data approval. Developers can contest ORCHA findings, and apps are anonymised \pc{zych2026}{4}{This secondary data analysis study gained ethical approval}.

\cf{Key findings} Median ORCHA score 60, below the threshold of 65 that ORCHA and the NHS consider adequate for public recommendation \pc{zych2026}{5}{ORCHA threshold score of 65 indicates}. No difference by cost type. Only 6.2\,\% of apps made the \emph{risks associated with using the app} clear. 97.3\,\% had an available privacy policy; 43\,\% stated GDPR compliance; 74\,\% informed users of their rights (Table 8) \pc{zych2026}{12}{Does the developer make clear risks associated with using the app}. Most common data types were usage data and email. Two clusters emerged, and the cluster collecting more data scored \emph{higher} \pc{zych2026}{14}{that mental health apps that collected more data, on average}.

\cf{Limitations stated} The tool's evolution and stringency, characteristics limited to OBR items, an exploratory design, and dependence on the instrument \pc{zych2026}{15}{This study is limited to the characteristics that were}. The dataset is not public.

\cf{Use for this thesis} (1) A contrast to a self-drawn store sample. Using a curated assessment dataset changes the population (downloads-ordered, ORCHA-assessed apps) and its representativeness. (2) The OBR items show that ``does the developer make clear risks associated with using the app'' is a recognised assessment criterion, met by only 6\,\%. This supports the thesis' focus on harm communication. (3) The finding that ``more data, higher score'' warns against treating compliance-style scores as privacy-protective. (4) Payment/cost types as descriptive variables (free, paid, in-app purchases).
\end{papercard}
